% ============================================================
%  BO_ML4F — Implementation Notes
%  Bayesian Optimisation with Gaussian Processes,
%  Rank-1 Cholesky Updates, Multi-Restart HPO,
%  Matern 5/2 Kernel, and Output Normalisation
%  Authors: Lecomte, Mendez
%  Last updated: July 2026
% ============================================================
\documentclass[11pt,a4paper]{article}

\usepackage[margin=2.5cm]{geometry}
\usepackage{amsmath,amssymb,amsthm}
\usepackage{bm}
\usepackage{booktabs}
\usepackage{hyperref}
\usepackage{xcolor}
\usepackage{listings}
\usepackage{algorithm}
\usepackage{algpseudocode}
\usepackage{graphicx}
\usepackage{microtype}

% ---------- Python listing style ----------
\definecolor{codegray}{rgb}{0.95,0.95,0.95}
\definecolor{keywordblue}{rgb}{0.0,0.3,0.7}
\definecolor{commentgreen}{rgb}{0.2,0.5,0.2}
\definecolor{stringred}{rgb}{0.6,0.1,0.1}

\lstdefinestyle{python}{
  language=Python,
  backgroundcolor=\color{codegray},
  basicstyle=\ttfamily\small,
  keywordstyle=\color{keywordblue}\bfseries,
  commentstyle=\color{commentgreen}\itshape,
  stringstyle=\color{stringred},
  numbers=left,
  numberstyle=\tiny\color{gray},
  numbersep=8pt,
  breaklines=true,
  frame=single,
  framesep=4pt,
  showstringspaces=false,
  tabsize=4,
}
\lstset{style=python}

% ---------- Notation ----------
\newcommand{\bx}{\bm{x}}
\newcommand{\by}{\bm{y}}
\newcommand{\bX}{\bm{X}}
\newcommand{\bK}{\bm{K}}
\newcommand{\bL}{\bm{L}}
\newcommand{\balpha}{\bm{\alpha}}
\newcommand{\btheta}{\bm{\theta}}
\newcommand{\R}{\mathbb{R}}
\newcommand{\N}{\mathcal{N}}
\newcommand{\GP}{\mathcal{GP}}
\newcommand{\EI}{\mathrm{EI}}
\newcommand{\PI}{\mathrm{PI}}
\newcommand{\LCB}{\mathrm{LCB}}
\newcommand{\tr}{\mathrm{tr}}
\newcommand{\diag}{\mathrm{diag}}

\theoremstyle{remark}
\newtheorem{remark}{Remark}

% ---------- File / line reference helpers ----------
\newcommand{\codefile}[1]{\texttt{#1}}
\newcommand{\codeline}[1]{line~#1}
\newcommand{\coderange}[2]{lines~#1--#2}
\newcommand{\codefn}[1]{\texttt{#1()}}

% ============================================================
\begin{document}

\title{\textbf{BO\_ML4F --- Implementation Notes}\\[4pt]
       \large Bayesian Optimisation with Gaussian Processes,\\
       Rank-1 Cholesky Updates, Multi-Restart HPO,\\
       Mat\'{e}rn 5/2 Kernel, and Output Normalisation}
\author{Y.\ Lecomte \and M.\ Mendez}
\date{July 2026}
\maketitle
\tableofcontents
\bigskip

% ============================================================
\section{Overview}
% ============================================================

This document describes the mathematics and Python implementation of
\codefile{BO\_ML4F.py}, a self-contained Bayesian Optimisation (BO) library
developed for the RSPA paper on the Burger equation control case.
A secondary helper module, \codefile{BO\_func\_YL.py}, provides a lower-level
implementation used in earlier prototypes.

The design goals are:
\begin{enumerate}
  \item \textbf{Correctness} --- exact GP posterior, proper acquisition functions
        for both minimisation and maximisation.
  \item \textbf{Efficiency} --- rank-1 Cholesky updates (Section~\ref{sec:rank1})
        reduce the per-iteration cost from $\mathcal{O}(n^3)$ to
        $\mathcal{O}(n^2)$ whenever the hyperparameters are held fixed.
  \item \textbf{Robustness} --- multi-restart hyperparameter optimisation with
        warm starting (Section~\ref{sec:hpo}) and output normalisation
        (Section~\ref{sec:normalisation}) for well-conditioned GP fitting.
  \item \textbf{Flexibility} --- two interchangeable kernels (squared-exponential
        and Mat\'{e}rn~5/2, Section~\ref{sec:kernels}) selectable via a single
        configuration field.
  \item \textbf{Modularity} --- all settings are grouped into typed dataclasses
        (\codefile{GPConfig}, \codefile{AcqConfig}, \codefile{OptimConfig},
        \codefile{BOConfig}, \codefile{SaveConfig}).
\end{enumerate}

\begin{remark}
All line numbers cited in this document refer to \codefile{BO\_ML4F.py} as it
stands after the July 2026 revision that introduced the rank-1 update,
multi-restart HPO, Mat\'{e}rn~5/2 kernel, and output normalisation.
\end{remark}

% ============================================================
\section{Kernel Functions}
\label{sec:kernels}
% ============================================================

The covariance function $k(\bx,\bx';\btheta)$ encodes prior beliefs about the
smoothness of the objective.  Two kernels are implemented, both parameterised
by the same triple $\btheta = (\ell, \sigma_f, \sigma_y)$ (length-scale,
signal amplitude, noise standard deviation).

% -----------------------------------
\subsection{Squared-Exponential (RBF) Kernel}
\label{sec:rbf}
% -----------------------------------

The squared-exponential kernel (also called the radial basis function, RBF,
or Gaussian kernel) is
\begin{equation}
  \label{eq:rbf}
  k_{\mathrm{SE}}(\bx,\bx') = \sigma_f^2
      \exp\!\Bigl(-\frac{\|\bx-\bx'\|^2}{2\ell^2}\Bigr)
    = \sigma_f^2 \exp\!\bigl(-\gamma\,\|\bx-\bx'\|^2\bigr),
  \quad \gamma = \frac{1}{2\ell^2}.
\end{equation}
This kernel is infinitely differentiable, implying infinitely smooth sample
paths.  It is a good default for well-behaved analytic functions but can
over-smooth physical simulation outputs.

\paragraph{Code.}
\codefn{rbf\_kernel\_amp} (\coderange{333}{341}) delegates pairwise
distance computation to \codefn{rbf\_kernel\_} in \codefile{BO\_func\_YL.py}
(\coderange{24}{26}):

\begin{lstlisting}[firstnumber=333, caption=\codefn{rbf\_kernel\_amp} in \codefile{BO\_ML4F.py}]
def rbf_kernel_amp(X1: Array, X2: Array, l_c: float, sigma_f: float) -> Array:
    gamma = 0.5 / (l_c**2)
    return (sigma_f**2) * rbf_kernel_(X1, X2, gamma=gamma)
\end{lstlisting}

\begin{lstlisting}[firstnumber=24, caption=\codefn{rbf\_kernel\_} in \codefile{BO\_func\_YL.py}]
def rbf_kernel_(X1, X2, gamma):
    sqdist = np.sum((X1[:, None, :] - X2[None, :, :])**2, axis=2)
    return np.exp(- gamma * sqdist)
\end{lstlisting}

The broadcasting trick \texttt{X1[:,None,:] - X2[None,:,:]} computes all
$n_1 \times n_2$ pairwise differences in one vectorised operation, avoiding
any Python loop.

% -----------------------------------
\subsection{Mat\'{e}rn 5/2 Kernel}
\label{sec:matern}
% -----------------------------------

The Mat\'{e}rn kernel with smoothness parameter $\nu = 5/2$ is
\begin{equation}
  \label{eq:matern52}
  k_{5/2}(\bx,\bx') = \sigma_f^2
  \left(1 + \frac{\sqrt{5}\,r}{\ell} + \frac{5\,r^2}{3\ell^2}\right)
  \exp\!\left(-\frac{\sqrt{5}\,r}{\ell}\right),
  \qquad r = \|\bx - \bx'\|_2,
\end{equation}
where $\ell$ and $\sigma_f$ play the same roles as in the RBF kernel.
Sample paths are exactly \emph{twice} differentiable --- a more realistic
prior for physical simulations whose outputs are not infinitely smooth.

For both kernels, evaluating at $r = 0$ gives $k(\bx,\bx) = \sigma_f^2$
(the bracketed factor in Eq.~\eqref{eq:matern52} equals unity when $r=0$).
This identity is exploited in Section~\ref{sec:predict} to avoid forming the
full $m \times m$ diagonal of $\bm{K}_{**}$.

\paragraph{Comparison with RBF.}
For a fixed length-scale $\ell$, the Mat\'{e}rn~5/2 kernel decays more
slowly with distance than the RBF kernel; equivalently, it places more prior
mass on non-smooth functions.  In practice this makes HPO more robust when
the objective has sharp features, and it is the default choice in production
libraries such as scikit-optimize.

\paragraph{Code.}
\codefn{matern52\_kernel\_amp} (\coderange{343}{353}):

\begin{lstlisting}[firstnumber=343, caption=\codefn{matern52\_kernel\_amp} in \codefile{BO\_ML4F.py}]
def matern52_kernel_amp(X1: Array, X2: Array, l_c: float, sigma_f: float) -> Array:
    sqdist = np.sum((X1[:, None, :] - X2[None, :, :])**2, axis=2)
    r      = np.sqrt(np.maximum(sqdist, 0.0))   # numerical safety for r~0
    a      = np.sqrt(5.0) * r / l_c
    return (sigma_f**2) * (1.0 + a + a**2 / 3.0) * np.exp(-a)
\end{lstlisting}

The \texttt{np.maximum(sqdist, 0.0)} guard prevents negative arguments to
\texttt{sqrt} that can arise from floating-point cancellation when two points
are nearly coincident.

% -----------------------------------
\subsection{Kernel Dispatcher}
\label{sec:dispatcher}
% -----------------------------------

All internal functions --- \codefn{gp\_fit}, \codefn{gp\_predict},
\codefn{negative\_log\_marginal\_likelihood}, \codefn{update\_gp\_rank1}
--- call a single dispatcher \codefn{kernel\_amp} (\coderange{356}{370})
that routes to the appropriate implementation based on a string key:

\begin{lstlisting}[firstnumber=356, caption=\codefn{kernel\_amp} dispatcher in \codefile{BO\_ML4F.py}]
def kernel_amp(X1: Array, X2: Array, l_c: float, sigma_f: float,
               kind: str = "rbf") -> Array:
    if kind == "rbf":
        return rbf_kernel_amp(X1, X2, l_c, sigma_f)
    elif kind == "matern52":
        return matern52_kernel_amp(X1, X2, l_c, sigma_f)
    else:
        raise ValueError(f"Unknown kernel '{kind}'. Choose 'rbf' or 'matern52'.")
\end{lstlisting}

The kernel choice is stored in \texttt{GPConfig.kernel} (\codeline{71}) and
propagated to \texttt{GPModel.kernel} (\codeline{136}) via
\codefn{build\_gp\_model} (\codeline{704}).  Switching kernels requires only
\texttt{GPConfig(kernel="matern52")} --- no other changes are needed.

% ============================================================
\section{Gaussian Process Regression}
\label{sec:gp}
% ============================================================

A Gaussian process is a distribution over functions:
\[
  f(\bx) \sim \GP\!\bigl(0,\, k(\bx,\bx';\btheta)\bigr).
\]
Given $n$ noisy observations $\mathcal{D}_n = \{(\bx_i, y_i)\}_{i=1}^n$
with $y_i = f(\bx_i)+\varepsilon_i$, $\varepsilon_i\sim\N(0,\sigma_y^2)$,
the posterior at a test point $\bx_*$ is
\begin{align}
  \mu_*      &= \bm{k}_*^\top \bigl(\bK_n + \sigma_y^2 \bm{I}\bigr)^{-1} \by_n, \\
  \sigma_*^2 &= k(\bx_*,\bx_*) - \bm{k}_*^\top \bigl(\bK_n + \sigma_y^2 \bm{I}\bigr)^{-1} \bm{k}_*,
\end{align}
where $[\bK_n]_{ij} = k(\bx_i,\bx_j)$ and
$\bm{k}_* = [k(\bx_i,\bx_*)]_{i=1}^n$.

% -----------------------------------
\subsection{GP Fitting via Cholesky Factorisation}
\label{sec:gpfit}
% -----------------------------------

Direct inversion of $\bm{K}_y = \bK_n + \sigma_y^2\bm{I}$ is $\mathcal{O}(n^3)$
and numerically fragile.  Instead we compute the Cholesky factorisation
\begin{equation}
  \bm{K}_y = \bL\bL^\top,
  \quad \bL \in \R^{n\times n}\ \text{lower triangular},
\end{equation}
and solve for the weight vector $\balpha$ in two triangular sweeps:
\begin{equation}
  \balpha = \bm{K}_y^{-1}\by_n
          = \bL^{-\top}\!\underbrace{\bL^{-1}\by_n}_{\bm{v}}.
\end{equation}
The posterior mean then reduces to the dot product $\mu_* = \bm{k}_*^\top\balpha$,
which reuses the precomputed $\balpha$ for all test points.

\paragraph{Code.}
\codefn{gp\_fit} (\coderange{373}{387}) assembles $\bm{K}_y$ using the
kernel dispatcher of Section~\ref{sec:dispatcher}, calls
\texttt{scipy.linalg.cholesky}, and solves for $\balpha$ via
\texttt{cho\_solve}:

\begin{lstlisting}[firstnumber=373, caption=\codefn{gp\_fit} in \codefile{BO\_ML4F.py}]
def gp_fit(Xs, ys, l_c=0.3, sigma_f=1.0, sigma_y=0.1, jitter=1e-10, kernel="rbf"):
    Xs = np.asarray(Xs, dtype=float)
    ys = np.asarray(ys, dtype=float).reshape(-1)

    Kss = kernel_amp(Xs, Xs, l_c=l_c, sigma_f=sigma_f, kind=kernel)
    n = Xs.shape[0]
    Ky = Kss + (sigma_y**2 + jitter) * np.eye(n)

    L = cholesky(Ky, lower=True)
    alpha = cho_solve((L, True), ys)
    return alpha, L
\end{lstlisting}

The \texttt{jitter} (default $10^{-10}$, \codeline{49}) is added to the
diagonal to ensure strict positive definiteness despite rounding errors.

% -----------------------------------
\subsection{GP Prediction}
\label{sec:predict}
% -----------------------------------

Given $(\balpha, \bL)$, prediction at $m$ test points $\bX_*\in\R^{m\times d}$
costs only $\mathcal{O}(mn)$:
\begin{align}
  \bm{\mu}_*   &= \bm{K}_{*n}\,\balpha, \\
  \bm{\sigma}_*^2 &= \diag(\bm{K}_{**}) - \diag\!\bigl(\bm{K}_{*n}\,\bm{V}\bigr),
  \quad \bm{V} = \bm{K}_y^{-1}\bm{K}_{n*}.
\end{align}
Since $k(\bx,\bx) = \sigma_f^2$ for both the RBF and Mat\'{e}rn 5/2 kernels
at $r=0$, we have $\diag(\bm{K}_{**}) = \sigma_f^2\,\bm{1}_m$, computed
cheaply as a scalar broadcast --- no full $m\times m$ matrix is needed.

\paragraph{Code.}
\codefn{gp\_predict} (\coderange{390}{413}):

\begin{lstlisting}[firstnumber=390, caption=\codefn{gp\_predict} in \codefile{BO\_ML4F.py}]
def gp_predict(X, Xs, alpha, L, l_c=0.3, sigma_f=1.0, return_cov=False, kernel="rbf"):
    Ks = kernel_amp(X, Xs, l_c=l_c, sigma_f=sigma_f, kind=kernel)
    mu = Ks @ alpha

    V = cho_solve((L, True), Ks.T)

    if return_cov:
        Kxx = kernel_amp(X, X, l_c=l_c, sigma_f=sigma_f, kind=kernel)
        cov = Kxx - Ks @ V
        return mu, cov
    else:
        # k(x,x) = sigma_f^2 for both RBF and Matern52
        Kxx_diag = (sigma_f**2) * np.ones(X.shape[0])
        var = Kxx_diag - np.sum(Ks * V.T, axis=1)
        return mu, var
\end{lstlisting}

When output normalisation is active, $\mu_*$ and $\sigma_*^2$ are returned
in the normalised space; denormalisation is the responsibility of the caller
(see Section~\ref{sec:normalisation}).

% -----------------------------------
\subsection{Output Normalisation}
\label{sec:normalisation}
% -----------------------------------

When the objective values $y$ have large absolute magnitude or units far from
unity, the GP hyperparameters $\sigma_f$ and $\sigma_y$ must also be large to
explain the data, which can make the NLL surface ill-conditioned and HPO slow
to converge.

The remedy is to work in a normalised space.  Before fitting, the observations
are centred and scaled:
\begin{equation}
  \label{eq:normalise}
  \tilde{y}_i = \frac{y_i - \bar{y}}{s_y}, \qquad
  \bar{y} = \frac{1}{n}\sum_{i=1}^n y_i, \quad
  s_y = \sqrt{\frac{1}{n}\sum_{i=1}^n (y_i - \bar{y})^2}.
\end{equation}
The GP is fitted on $\tilde{\by}$; the resulting posterior mean
$\tilde\mu_*$ and variance $\tilde\sigma_*^2$ are denormalised at call sites:
\begin{equation}
  \label{eq:denormalise}
  \mu_* = \tilde\mu_* \cdot s_y + \bar{y},
  \qquad
  \sigma_*^2 = \tilde\sigma_*^2 \cdot s_y^2.
\end{equation}
This keeps $\sigma_f \sim 1$ and $\sigma_y \sim 1$ regardless of the
objective scale, and all acquisition functions (EI, PI, LCB) operate in the
original scale after denormalisation.

\begin{remark}
The normalisation statistics $(\bar{y}, s_y)$ are computed once per full refit
and stored in \texttt{gp.y\_mean} and \texttt{gp.y\_std}
(\coderange{139}{140} of \codefile{BO\_ML4F.py}).  Rank-1 updates reuse the
existing statistics --- they do \emph{not} recompute $\bar{y}$ and $s_y$,
because updating those would invalidate the Cholesky factor built under the
previous normalisation.
\end{remark}

\paragraph{Code.}
The normalisation block is inside \codefn{fit\_gp} (\coderange{715}{760}),
applied before HPO and before the Cholesky call:

\begin{lstlisting}[firstnumber=726, caption=Output normalisation in \codefn{fit\_gp}]
    if gp_cfg.normalize_y:
        y_mean = float(np.mean(y_raw))
        y_std  = float(np.std(y_raw))
        if y_std < 1e-12:
            y_std = 1.0          # guard against constant outputs
        gp.y_mean = y_mean
        gp.y_std  = y_std
        gp.ys = (y_raw - y_mean) / y_std
    else:
        gp.y_mean = 0.0
        gp.y_std  = 1.0
        gp.ys = y_raw
\end{lstlisting}

Denormalisation [Eq.~\eqref{eq:denormalise}] is applied in
\codefn{make\_acquisition} immediately after the \codefn{gp\_predict} call:

\begin{lstlisting}[firstnumber=784, caption=Denormalisation in \codefn{make\_acquisition}]
        mu_norm, var_norm = gp_predict(
            Xcand, gp.Xs, gp.alpha, gp.L,
            l_c=gp.l_c, sigma_f=gp.sigma_f,
            return_cov=False, kernel=gp.kernel,
        )
        # Eq. (4): denormalise to original y scale
        mu  = mu_norm  * gp.y_std + gp.y_mean
        var = var_norm * gp.y_std**2
\end{lstlisting}

The same denormalisation is applied in \codefn{plt\_state} (\codeline{957})
so that plotted GP posteriors are always in physical units.

% ============================================================
\section{Hyperparameter Optimisation (HPO)}
\label{sec:hpo}
% ============================================================

% -----------------------------------
\subsection{Log Marginal Likelihood}
% -----------------------------------

The hyperparameters $\btheta = (\ell, \sigma_f, \sigma_y)$ are tuned by
maximising the log marginal likelihood (equivalently, minimising the negative
log marginal likelihood, NLL):
\begin{equation}
  \label{eq:lml}
  \log p(\tilde{\by} \mid \bX, \btheta)
  = -\frac{1}{2}\tilde{\by}^\top \bm{K}_y^{-1}\tilde{\by}
    - \sum_{i=1}^n \log L_{ii}
    - \frac{n}{2}\log(2\pi),
\end{equation}
where $\tilde{\by}$ denotes the (possibly normalised) observations stored
in \texttt{gp.ys}.  Because all three hyperparameters are strictly positive,
optimisation is performed in log-space $\tilde\btheta = \log\btheta$, so
L-BFGS-B operates over $\R^3$ with optional box constraints
(\texttt{GPConfig.theta\_bounds\_log}, \codeline{58}).

\paragraph{Code.}
\codefn{negative\_log\_marginal\_likelihood} (\coderange{568}{602}) evaluates
Eq.~\eqref{eq:lml}.  A \texttt{kernel} argument ensures the NLL is computed
with the correct covariance function:

\begin{lstlisting}[firstnumber=568, caption=NLL in \codefile{BO\_ML4F.py}]
def negative_log_marginal_likelihood(
    theta_log, X_train, y_train,
    jitter=1e-10, MLE_hist=None, param_hist=None,
    kernel="rbf",
):
    l_c, sigma_f, sigma_y = np.exp(theta_log)
    K  = kernel_amp(X_train, X_train, l_c=l_c, sigma_f=sigma_f, kind=kernel)
    Ky = K + (sigma_y**2 + jitter) * np.eye(X_train.shape[0])
    L  = cholesky(Ky, lower=True)
    alpha = cho_solve((L, True), y_train)
    ll = -0.5*(y_train @ alpha) - np.sum(np.log(np.diag(L))) \
         - 0.5*len(y_train)*np.log(2.0*np.pi)
    return -float(ll)
\end{lstlisting}

% -----------------------------------
\subsection{Multi-Restart Strategy with Warm Starting}
\label{sec:multistart}
% -----------------------------------

The NLL surface can have multiple local optima, especially when data are sparse.
\codefn{optimize\_gp\_hyperparams} (\coderange{605}{697}) addresses this with
a three-tier starting-point strategy:

\begin{enumerate}
  \item \textbf{Warm start} (\codeline{640}): the solution $\tilde\btheta^*$
        from the \emph{previous} BO iteration is reused.  Adding one point
        rarely shifts the NLL optimum much, so this typically converges in
        very few L-BFGS-B steps.
  \item \textbf{User or default start} (\coderange{644}{647}): either
        \texttt{gp\_cfg.theta0\_log} or the current hyperparameters in
        log-space.
  \item \textbf{Random restarts} (\coderange{650}{652}):
        \texttt{n\_hpo\_restarts} (default 3, \codeline{65}) vectors drawn
        from $\N(\bm{0},\bm{I})$ provide global coverage to escape local optima.
\end{enumerate}

The restart with the lowest NLL is kept (\coderange{657}{670}) and its
solution stored in \texttt{gp.theta\_log} (\codeline{695}) to seed the next
warm start.

\begin{lstlisting}[firstnumber=637, caption=Multi-restart loop in \codefn{optimize\_gp\_hyperparams}]
    x0_list: List[Array] = []

    if gp.theta_log is not None:          # (1) warm start
        x0_list.append(gp.theta_log.copy())
    if gp_cfg.theta0_log is not None:     # (2) user / default
        x0_list.append(np.asarray(gp_cfg.theta0_log, dtype=float).reshape(-1))
    else:
        x0_list.append(np.log([gp.l_c, gp.sigma_f, gp.sigma_y]))
    rng_hpo = np.random.default_rng()    # (3) random restarts
    for _ in range(max(0, gp_cfg.n_hpo_restarts)):
        x0_list.append(rng_hpo.standard_normal(3))

    best_nll, best_theta = np.inf, None
    for x0 in x0_list:
        res = minimize(negative_log_marginal_likelihood, x0=x0,
                       args=(X, y, gp.jitter, None, None, gp.kernel),
                       method="L-BFGS-B", bounds=bounds)
        if np.isfinite(res.fun) and res.fun < best_nll:
            best_nll, best_theta = res.fun, res.x.copy()
    ...
    gp.theta_log = best_theta.copy()
\end{lstlisting}

\paragraph{HPO scheduling.}
\texttt{gp\_cfg.hpo\_every} (default~1, \codeline{53}) controls the period:
\texttt{hpo\_every=5} triggers HPO only every fifth iteration, substantially
reducing overhead for cheap objectives.

% ============================================================
\section{Acquisition Functions}
\label{sec:acq}
% ============================================================

Acquisition functions quantify the utility of sampling at a candidate point
$\bx$ given the GP posterior $(\mu_*, \sigma_*)$ in the \emph{original}
(denormalised) scale.  All three functions are returned by
\codefn{make\_acquisition} (\coderange{766}{838}) as a vectorised callable
$a: \R^{m\times d} \to \R^m$.

\subsection{Expected Improvement (EI)}

For \emph{minimisation} with current best $y^\dagger = \min_i y_i$ and
jitter $\xi \geq 0$:
\begin{equation}
  \label{eq:ei}
  \EI(\bx) = (y^\dagger - \mu_* - \xi)\,\Phi(Z)
             + \sigma_*\,\phi(Z),
  \quad Z = \frac{y^\dagger - \mu_* - \xi}{\sigma_*},
\end{equation}
where $\Phi$ and $\phi$ are the standard-normal CDF and PDF.  EI is zero
wherever $\sigma_* = 0$ (no uncertainty, no improvement possible).

\subsection{Probability of Improvement (PI)}
\begin{equation}
  \PI(\bx) = \Phi\!\Bigl(\frac{y^\dagger - \mu_* - \xi}{\sigma_*}\Bigr).
\end{equation}
PI ignores the \emph{magnitude} of improvement and is more exploitative
than EI.

\subsection{Lower Confidence Bound (LCB/UCB)}

For minimisation the acquisition is
\begin{equation}
  \LCB(\bx) = \mu_* - \kappa\,\sigma_*,
\end{equation}
which the internal optimiser maximises as $-\LCB = -\mu_* + \kappa\sigma_*$.
The parameter $\kappa$ (\codeline{73}) trades exploration against exploitation.

\begin{remark}
\textbf{Sign convention.}  The optimiser always \emph{maximises} $a(\bx)$.
For EI and PI this is natural; for LCB the code returns
$-(\mu - \kappa\sigma)$ so that maximising it is equivalent to minimising the
lower confidence bound.
\end{remark}

% ============================================================
\section{Acquisition Optimisation}
\label{sec:acqopt}
% ============================================================

Finding $\bx_{\mathrm{next}} = \arg\max_{\bx\in\mathcal{X}} a(\bx)$ is
non-convex.  Three strategies are available in \codefn{optimize\_acquisition}
(\coderange{856}{951}), selected by \texttt{OptimConfig.method}.

\paragraph{Random} (\texttt{method="random"}): draw $N$ candidates uniformly
from the bounding box (\texttt{n\_raw\_samples}, default 2000) and return
$\arg\max$.

\paragraph{Grid} (\texttt{method="grid"}): build a uniform Cartesian grid with
\texttt{grid\_n\_per\_dim} points per axis (default 20).  Suitable only for
$d \leq 3$ since the grid grows exponentially with dimension.

\paragraph{Refined} (\texttt{method="refined"}): combine the global search
above with local L-BFGS-B polishing.  The top \texttt{n\_restarts} (default 10)
candidates seed \texttt{scipy.optimize.minimize} for precise local convergence.
Recommended when the objective is expensive and precision matters.

\begin{remark}
Because $a$ reuses the precomputed $(\balpha, \bL)$, each evaluation at $m$
points costs $\mathcal{O}(mn)$.  The random global scan costs $\mathcal{O}(Nn)$
total; the refined step adds $\mathcal{O}(K T n)$ where $K$ is the number
of restarts and $T$ the average L-BFGS-B step count.
\end{remark}

% ============================================================
\section{Rank-1 Cholesky Update}
\label{sec:rank1}
% ============================================================

% -----------------------------------
\subsection{Motivation and Complexity}
% -----------------------------------

At each BO iteration the GP is refitted on a dataset growing by one point.
The standard approach assembles the full $n\times n$ kernel matrix and calls
\texttt{scipy.linalg.cholesky} at cost $\mathcal{O}(n^3)$.  After $T$
iterations with $n_0$ initial points, total fitting cost is
$\mathcal{O}\bigl((n_0+T)^4\bigr)$.  The rank-1 update reduces each
per-iteration cost to $\mathcal{O}(n^2)$, giving
$\mathcal{O}\bigl((n_0+T)^3\bigr)$ total --- an asymptotic factor of
$(n_0+T)$ improvement.

% -----------------------------------
\subsection{Block Cholesky Extension}
% -----------------------------------

Given $\bm{K}_n^y = \bK_n + \sigma_y^2\bm{I}_n = \bL_n\bL_n^\top$, after
adding $\bx_{n+1}$ the augmented matrix is
\begin{equation}
  \bm{K}_{n+1}^y =
  \begin{pmatrix}
    \bm{K}_n^y & \bm{k} \\
    \bm{k}^\top & k_{n+1,n+1} + \sigma_y^2 + \delta
  \end{pmatrix},
\end{equation}
where $\bm{k} = [k(\bx_i, \bx_{n+1})]_{i=1}^n$ and $\delta$ is the jitter.
The new Cholesky factor inherits the block structure
\begin{equation}
  \label{eq:chol_block}
  \bL_{n+1} =
  \begin{pmatrix}
    \bL_n & \bm{0} \\
    \bm{l}^\top & l_{nn}
  \end{pmatrix},
\end{equation}
with the two new entries determined by matching
$\bL_{n+1}\bL_{n+1}^\top = \bm{K}_{n+1}^y$:
\begin{align}
  \label{eq:chol_l}
  \bm{l}   &= \bL_n^{-1}\bm{k}
             \quad \text{(forward substitution, $\mathcal{O}(n^2)$)}, \\
  \label{eq:chol_lnn}
  l_{nn}   &= \sqrt{k_{n+1,n+1} + \sigma_y^2 + \delta - \bm{l}^\top\bm{l}}.
\end{align}

\paragraph{Code.} \codefn{chol\_rank1\_update} (\coderange{420}{465}):

\begin{lstlisting}[firstnumber=454, caption=Cholesky extension in \codefn{chol\_rank1\_update}]
    n = L.shape[0]
    l = solve_triangular(L, k_cross, lower=True)   # Eq. (7): O(n^2)
    d    = k_self + sigma_y**2 + jitter - float(np.dot(l, l))
    l_nn = np.sqrt(max(d, jitter))                 # Eq. (8): guard for r~0
    L_new = np.zeros((n+1, n+1), dtype=float)
    L_new[:n, :n] = L
    L_new[n, :n]  = l
    L_new[n, n]   = l_nn
    return L_new
\end{lstlisting}

% -----------------------------------
\subsection{Alpha Update}
% -----------------------------------

After the Cholesky extension we need $\balpha_{n+1} = \bm{K}_{n+1}^{y,-1}\tilde{\by}_{n+1}$,
where all $y$ values are in normalised space when \texttt{normalize\_y=True}.
Using the block structure of $\bL_{n+1}$ the forward and backward sweeps give:
\begin{equation}
  \label{eq:alpha_update}
  \boxed{
  \bm{v}_1 = \bL_n^\top\balpha_n, \quad
  v_2 = \frac{\tilde{y}_{n+1}-\bm{l}^\top\bm{v}_1}{l_{nn}}, \quad
  \alpha_{n+1}^{(n+1)} = \frac{v_2}{l_{nn}}, \quad
  \balpha_{n+1}^{(1:n)} = \bL_n^{-\top}\!\bigl(\bm{v}_1-\bm{l}\,\alpha_{n+1}^{(n+1)}\bigr).
  }
\end{equation}
The key insight is that $\bm{v}_1 = \bL_n^{-1}\tilde{\by}_n = \bL_n^\top\balpha_n$
(since $\balpha_n = \bL_n^{-\top}\bL_n^{-1}\tilde{\by}_n$), so $\bm{v}_1$
is recovered by a cheap $\mathcal{O}(n^2)$ mat-vec, without an extra
triangular solve.  Total cost: one mat-vec plus one backward solve, both
$\mathcal{O}(n^2)$.

\paragraph{Code.} \codefn{alpha\_rank1\_update} (\coderange{468}{512}):

\begin{lstlisting}[firstnumber=497, caption=Alpha update in \codefn{alpha\_rank1\_update}]
    l    = L_new[n, :n]
    l_nn = L_new[n, n]

    # v_1 = L_old^T alpha_old  [Eq. (9), O(n^2) mat-vec, no extra solve]
    v1 = L_old.T @ alpha_old

    v2             = (y_new - float(np.dot(l, v1))) / l_nn
    alpha_new_last = v2 / l_nn

    rhs             = v1 - l * alpha_new_last
    alpha_new_first = solve_triangular(L_old, rhs, lower=True, trans='T')

    return np.concatenate([alpha_new_first, [alpha_new_last]])
\end{lstlisting}

% -----------------------------------
\subsection{Integration into the BO Loop}
% -----------------------------------

\codefn{update\_gp\_rank1} (\coderange{515}{561}) assembles the cross-covariance
$\bm{k}$ and uses $k_{n+1,n+1} = \sigma_f^2$ (valid for both kernels at $r=0$).
When normalisation is active the new observation is mapped to normalised space
using the \emph{existing} $(\bar{y}, s_y)$ before the update:

\begin{lstlisting}[firstnumber=539, caption=\codefn{update\_gp\_rank1}]
    y_new_norm = (y_new_raw - gp.y_mean) / gp.y_std  # existing stats: no recompute

    k_cross = kernel_amp(gp.Xs, x_new, gp.l_c, gp.sigma_f, kind=gp.kernel).reshape(-1)
    k_self  = float(gp.sigma_f**2)   # k(x,x) = sigma_f^2 for both kernels

    L_new     = chol_rank1_update(gp.L, k_cross, k_self, gp.sigma_y, gp.jitter)
    alpha_new = alpha_rank1_update(gp.L, gp.alpha, L_new, y_new_norm)
\end{lstlisting}

In the main BO loop (\coderange{1144}{1173}) three fitting paths are selected:

\begin{lstlisting}[firstnumber=1164, caption=Fitting path selection in \codefn{bayesian\_optimization}]
        if it == 0 or do_hpo:
            gp = fit_gp(gp, X, y, gp_cfg, it)          # (a) O(n^3), renormalise
        elif gp_cfg.use_rank1_update and gp.L is not None:
            gp = update_gp_rank1(gp, X[-1:], float(y[-1]))  # (b) O(n^2) rank-1
        else:
            gp = fit_gp(gp, X, y, gp_cfg, it)          # (c) full, HPO skipped
\end{lstlisting}

\begin{itemize}
  \item \textbf{Path (a)} is mandatory at iteration~0 and whenever HPO runs.
        New hyperparameters invalidate the entire kernel matrix, and the
        normalisation statistics are recomputed from scratch.
  \item \textbf{Path (b)} is triggered when \texttt{use\_rank1\_update=True}
        and no HPO occurs.  At this point \texttt{X} and \texttt{y} already
        contain the new point while \texttt{gp.Xs / gp.ys} still hold the
        previous dataset, so \texttt{X[-1:]} and \texttt{y[-1]} identify the
        single new observation.
  \item \textbf{Path (c)} is the fallback full refit with HPO skipped.
\end{itemize}

\begin{remark}
When rank-1 updates are interleaved with periodic HPO (e.g.\
\texttt{hpo\_every=5}), the pattern is: full refit at iterations
$0, 5, 10, \ldots$ (with renormalisation and HPO), followed by four
$\mathcal{O}(n^2)$ rank-1 steps using the frozen normalisation.  The periodic
full refit also resets any accumulated floating-point drift in the factors.
\end{remark}

% ============================================================
\section{The BO Loop}
\label{sec:bo_loop}
% ============================================================

The complete algorithm is in \codefn{bayesian\_optimization}
(\coderange{1089}{end}).  Algorithm~\ref{alg:bo} provides the pseudocode.

\begin{algorithm}[ht]
\caption{Bayesian Optimisation --- \codefn{bayesian\_optimization}}
\label{alg:bo}
\begin{algorithmic}[1]
\Require objective $f$, bounds $\mathcal{X}$, configs
  \texttt{bo\_cfg, gp\_cfg, acq\_cfg, optim\_cfg}
\State Sample $n_0$ initial points (uniform or LHS, \codeline{1128})
\State Build GP container $\leftarrow$ \codefn{build\_gp\_model} (\codeline{1139})
\For{$t = 0, 1, \ldots, T-1$}
  \If{$t = 0$ \textbf{or} HPO scheduled}
    \State Full refit incl.\ (re-)normalisation: $(\balpha, \bL) \leftarrow$ \codefn{fit\_gp}
           \hfill \textit{Path (a)}
  \ElsIf{\texttt{use\_rank1\_update}}
    \State $(\balpha, \bL) \leftarrow$ \codefn{update\_gp\_rank1}
           \hfill \textit{Path (b)}
  \Else
    \State Full refit, HPO skipped
           \hfill \textit{Path (c)}
  \EndIf
  \State $y^\dagger \leftarrow \min_i y_i$ \quad (raw scale)
  \State $a(\cdot) \leftarrow$ \codefn{make\_acquisition}$(acq\_cfg, gp, y^\dagger)$
         \quad [GP in norm.\ space; $a$ in raw scale]
  \State $\bx_{t+1} \leftarrow$ \codefn{optimize\_acquisition}$(a, \mathcal{X}, optim\_cfg)$
  \State $y_{t+1} \leftarrow f(\bx_{t+1})$
  \State Append $(\bx_{t+1}, y_{t+1})$ to $(X, y)$
\EndFor
\State \Return $(\bx^*, y^*) = \arg\min_{(\bx,y)\in\mathcal{D}} y$
\end{algorithmic}
\end{algorithm}

% ============================================================
\section{Configuration Reference}
\label{sec:config}
% ============================================================

All settings live in typed \texttt{dataclass} containers defined in
\coderange{43}{189} of \codefile{BO\_ML4F.py}.
Table~\ref{tab:config} summarises the most important parameters.

\begin{table}[ht]
\centering
\caption{Key configuration parameters.}
\label{tab:config}
\small
\begin{tabular}{llll}
\toprule
\textbf{Class} & \textbf{Field} & \textbf{Default} & \textbf{Description}\\
\midrule
\texttt{GPConfig} & \texttt{l\_c} & 0.3
  & Length-scale $\ell$ \\
 & \texttt{sigma\_f} & 1.0
  & Kernel amplitude $\sigma_f$ \\
 & \texttt{sigma\_y} & 0.1
  & Noise std $\sigma_y$ \\
 & \texttt{jitter} & $10^{-10}$
  & Diagonal stabiliser $\delta$ \\
 & \texttt{kernel} & \texttt{"rbf"}
  & \texttt{"rbf"} or \texttt{"matern52"} (\S\ref{sec:kernels}) \\
 & \texttt{normalize\_y} & \texttt{False}
  & Centre/scale $y$ before fitting (\S\ref{sec:normalisation}) \\
 & \texttt{optimize\_hyperparams} & \texttt{False}
  & Enable HPO \\
 & \texttt{hpo\_every} & 1
  & Run HPO every $k$ iterations \\
 & \texttt{n\_hpo\_restarts} & 3
  & Random restarts in HPO \\
 & \texttt{use\_rank1\_update} & \texttt{False}
  & Enable $\mathcal{O}(n^2)$ update (\S\ref{sec:rank1}) \\
\midrule
\texttt{AcqConfig} & \texttt{kind} & \texttt{"EI"}
  & Acquisition type (EI/PI/UCB) \\
 & \texttt{xi} & 0.01
  & Exploration jitter $\xi$ (EI/PI) \\
 & \texttt{kappa} & 2.0
  & Exploration weight $\kappa$ (LCB) \\
 & \texttt{maximize} & \texttt{False}
  & Set \texttt{True} for maximisation \\
\midrule
\texttt{OptimConfig} & \texttt{method} & \texttt{"random"}
  & Acq.\ optimisation mode \\
 & \texttt{n\_raw\_samples} & 2000
  & Candidates for random/refined \\
 & \texttt{n\_restarts} & 10
  & Local restarts in refined mode \\
\midrule
\texttt{BOConfig} & \texttt{n\_init} & 5
  & Initial points \\
 & \texttt{n\_iter} & 20
  & BO iterations \\
 & \texttt{init\_sampling} & \texttt{"uniform"}
  & Initial design (\texttt{"lhs"} available) \\
\bottomrule
\end{tabular}
\end{table}

\paragraph{Recommended configuration for a medium-cost problem.}

\begin{lstlisting}[numbers=none]
gp_cfg = bo.GPConfig(
    kernel               = "matern52",  # twice-differentiable, more robust
    normalize_y          = True,        # keeps sigma_f ~ 1 regardless of y scale
    optimize_hyperparams = True,
    hpo_every            = 5,           # full refit every 5 iterations
    n_hpo_restarts       = 3,           # warm + default + 3 random = 5 starts
    use_rank1_update     = True,        # O(n^2) between HPO calls
)
\end{lstlisting}

% ============================================================
\section{Test Cases and Examples}
\label{sec:examples}
% ============================================================

The folder contains five test scripts organised into three groups:
two examples for \codefile{BO\_ML4F.py}, two legacy examples for
\codefile{BO\_func\_YL.py}, and one reference script based on
scikit-optimize.

% -----------------------------------
\subsection{\texttt{1D\_Test\_CASE.py} --- 1D example with \codefile{BO\_ML4F.py}}
\label{sec:ex1d}
% -----------------------------------

This is the primary smoke-test for the new library.  The objective is the
one-dimensional noisy function
\begin{equation}
  \label{eq:f1d}
  f(x) = \sin(5x)\bigl(1 - \tanh(x^2)\bigr) + \varepsilon,
  \quad \varepsilon \sim \N(0, 0.1^2),
  \quad x \in [-2,\, 2].
\end{equation}
The function has a clear global minimum near $x \approx 0$ flanked by
shallower local minima; the tanh factor suppresses the oscillations away
from the origin.  The script first plots 100 dense evaluations for visual
inspection, then launches \codefn{bayesian\_optimization} with all default
settings:

\begin{lstlisting}[numbers=none, caption=\codefile{1D\_Test\_CASE.py} --- key call]
res = bo.bayesian_optimization(
    f=func,
    bounds=[(-2.0, 2.0)],
    bo_cfg=bo.BOConfig(random_state=1234),
    gp_cfg=bo.GPConfig(),          # RBF kernel, no HPO, no rank-1
    acq_cfg=bo.AcqConfig(),        # EI, xi=0.01
    optim_cfg=bo.OptimConfig(),    # random search, N=2000
    save_cfg=bo.SaveConfig(
        plt_state_enabled=True,    # plot GP + EI at every iteration
        save_path="./out"
    ),
)
\end{lstlisting}

Because $d=1$, \texttt{plt\_state\_enabled=True} is safe: the 1D plotting
routine in \codefn{plt\_state} draws the true function, the GP posterior
mean $\pm\,2\sigma$, the observations, and the acquisition function at every
BO iteration, saving each frame to \texttt{./out/}.

% -----------------------------------
\subsection{\texttt{3D\_Test\_CASE.py} --- 3D example with \codefile{BO\_ML4F.py}}
\label{sec:ex3d}
% -----------------------------------

This script tests the library on a three-dimensional objective:
\begin{equation}
  \label{eq:f3d}
  f(\bx) = \sin(3x_0) + \tfrac{1}{2}\cos(5x_1) + \tfrac{1}{5}x_2^2 + \varepsilon,
  \quad \varepsilon \sim \N(0, 0.1^2),
  \quad \bx \in [-2,\,2]^3.
\end{equation}
The script first generates 500 random points via \codefn{init\_dataset}
and plots them in a 3D scatter coloured by $f(\bx)$, providing a visual
sanity check of the objective's structure.  It then runs the BO loop:

\begin{lstlisting}[numbers=none, caption=\codefile{3D\_Test\_CASE.py} --- key call]
res = bo.bayesian_optimization(
    f=func,
    bounds=[(-2, 2), (-2, 2), (-2, 2)],
    bo_cfg=bo.BOConfig(
        n_init=10,
        n_iter=30,
        random_state=1234
    ),
    gp_cfg=bo.GPConfig(),
    acq_cfg=bo.AcqConfig(kind="EI"),
    optim_cfg=bo.OptimConfig(method="random"),
    save_cfg=bo.SaveConfig(
        plt_state_enabled=False,   # REQUIRED for d > 1
        save_path="./out_3d"
    ),
)
\end{lstlisting}

\texttt{plt\_state\_enabled} must be \texttt{False} for $d>1$: the plotting
routine only supports 1D inputs (it builds a dense grid on the single axis).
The best solution is accessed via \texttt{res.best\_x} and
\texttt{res.best\_y} after the loop.

% -----------------------------------
\subsection{\texttt{1D\_Test\_CASE\_YL.py} and \texttt{3D\_Test\_CASE\_YL.py}
  --- Legacy examples}
\label{sec:exyl}
% -----------------------------------

These two scripts use the lower-level \codefile{BO\_func\_YL.py} module and
do not call \codefn{bayesian\_optimization}.  They are kept as reference
implementations and pedagogical baselines.

\textbf{1D\_Test\_CASE\_YL.py} implements the complete BO loop manually for
the same 1D function (Eq.~\eqref{eq:f1d}).  It calls \codefn{log\_marginal\_likelihood},
\codefn{gp\_fit}, \codefn{gp\_predict}, and \codefn{propose\_location} directly,
making it easy to trace every step of the algorithm without abstraction layers.

\textbf{3D\_Test\_CASE\_YL.py} is an \emph{incomplete} script that generates
a 3D dataset with a Latin Hypercube sampler (\texttt{scipy.stats.qmc.LatinHypercube}),
plots the input cloud coloured by the objective
\[
  g(\bx) = \sin(x_0) + \tfrac{1}{2}\cos(x_1)
          + \tfrac{1}{4}x_2^2 + 0.3\sin(x_0 x_1),
\]
and fits a GP to the data --- but does not yet run a BO loop.  It serves as
a template for future extensions.

% -----------------------------------
\subsection{\texttt{3\_Bayesian\_Optimization\_Scipy\_Optimize.py}
  --- scikit-optimize reference}
\label{sec:exskopt}
% -----------------------------------

This script, adapted from the scikit-optimize tutorial, applies
\texttt{skopt.gp\_minimize} to the same 1D objective (Eq.~\eqref{eq:f1d}).
It was used as the reference when benchmarking features absent from the
original \codefile{BO\_func\_YL.py} implementation:

\begin{lstlisting}[numbers=none, caption=\codefile{3\_Bayesian\_Optimization\_Scipy\_Optimize.py}]
from skopt import gp_minimize
res = gp_minimize(f,
                  [(-2.0, 2.0)],
                  acq_func="EI",
                  n_calls=15,
                  n_random_starts=5,
                  noise=0.1**2,
                  random_state=1234)
\end{lstlisting}

The script also produces an animated GIF showing the GP posterior and EI
acquisition at each of the first ten iterations, using
\texttt{skopt.plots.plot\_gaussian\_process}.  The two features identified
in that comparison and subsequently ported into \codefile{BO\_ML4F.py} are
the Mat\'{e}rn~5/2 kernel (Section~\ref{sec:matern}) and output normalisation
(Section~\ref{sec:normalisation}).

% ============================================================
\section{Corrections to \texttt{BO\_func\_YL.py}}
\label{sec:funcyl}
% ============================================================

Two issues in the legacy module \codefile{BO\_func\_YL.py} were fixed in the
July 2026 revision.

\paragraph{Vectorised EI.}
The original \codefn{propose\_location} called \codefn{EI} in a Python
\texttt{for} loop over all grid points.  The revised version passes the entire
grid matrix \texttt{grid\_points} (shape $N_g \times d$) to a batch-capable
\codefn{EI} that calls the vectorised \codefn{gp\_predict}:
\begin{lstlisting}[numbers=none]
# Before (loop):
ei_values = np.array([EI(x, ...)[0] for x in grid_points])
# After (vectorised):
ei_values = EI(grid_points, X_train, y_train, params, xi, alpha, L)
\end{lstlisting}
The new \codefn{EI} (\coderange{85}{103} of \codefile{BO\_func\_YL.py})
accepts an $(N,d)$ batch, returns an $(N,)$ array, and correctly zeros
entries where $\sigma = 0$.

\paragraph{Stray tab character.}
A mixed-indent bug (tab + spaces) in \codefn{plot\_n\_save\_BO} caused an
\texttt{IndentationError} under strict Python parsers; replaced with four
spaces throughout.

\end{document}
